% =====================================================================
\section{Curriculum-learning study}
\label{sec:res-cur}

This appendix summarises the curriculum-learning extension of
the forward PINN; complete tables, figures and per-variant
extraction results are provided in the repository
documentation.

\subsection{Method}
\label{sec:cur-pinn}

The temporal domain is partitioned into $K$ windows trained
sequentially: each window is an independent network with the
forward baseline's architecture, loss weights and optimiser
schedule, and the trained $k$-th network supplies the initial
condition of the $(k{+}1)$-th window by evaluation at the split
time. The three configurations (Cur-2w, Cur-3w, Cur-soft;
Table~\ref{tab:cur-summary}) are motivated by the
single-window forward PINN concentrating $\sim\!41\%$ of its
RMSD in the late-time band $t \in [37.5,50]\,M$.

\subsection{Results}

Table~\ref{tab:cur-summary} reports full-domain and per-window
field accuracy against the FD reference.

\begin{table}[!htbp]
\centering
\caption{Curriculum PINN variants versus the single-window
forward baseline: field accuracy. Cur-2w is a hard two-window
partition at $t = 25\,M$; Cur-3w a hard three-window partition
at $t = 17, 34\,M$; Cur-soft a variant with a $10\,M$ overlap,
continuity-anchor loss and smoothstep blending. Columns:
full-domain RMSD against the FD reference and the per-window
RMSDs on each sub-domain (``---'' if the variant does not use
that window); for Cur-soft the third ``$W$'' column reports the
RMSD on the soft-blend overlap region $[20,30]\,M$ rather than
a third disjoint window.}
\label{tab:cur-summary}
\footnotesize
\setlength{\tabcolsep}{4pt}
\renewcommand{\arraystretch}{1.05}
\begin{tabular}{@{}lcccc@{}}
\toprule
Variant & Full-domain RMSD & $W_{1}$ RMSD & $W_{2}$ RMSD & $W_{3}$ / overlap RMSD \\
        & ($\times 10^{-3}$) & ($\times 10^{-3}$) & ($\times 10^{-3}$) & ($\times 10^{-3}$) \\
\midrule
Baseline (single window) & $\mathbf{0.817}$ & $-\!\!-\!\!-$ & $-\!\!-\!\!-$ & $-\!\!-\!\!-$ \\
\midrule
Cur-2w   & $0.845$ & $\mathbf{0.516}$ & $1.077$ & $-\!\!-\!\!-$ \\
Cur-3w   & $0.976$ & $0.554$          & $0.920$ & $1.322$ \\
Cur-soft & $3.155$ & $1.739$          & $3.932$ & $2.394$ (overlap) \\
\bottomrule
\end{tabular}
\end{table}

The two-window curriculum delivers the intended local gain: its
first window reaches RMSD $5.16\times 10^{-4}$, a $37\%$
improvement over the baseline restricted to the same
sub-domain. The gain does not survive stitching: the second
window inherits the first's residual error through the
numerical IC transfer and must additionally resolve the
late-time decay, ending at $1.08\times 10^{-3}$, so the
full-domain RMSD closes $3\%$ above the baseline; Cur-3w
compounds one interpolation error per splice
($0.55 \to 0.92 \to 1.32$, units of $10^{-3}$) and ends $19\%$
above; Cur-soft performs worst, at roughly four times the
baseline, with no M4 plateau at all. Extraction is insensitive
to the curriculum at the wide-window M2 level but degrades at
plateau level (Cur-2w: $\sigma_\omega = 2\times 10^{-4} \to
9\times 10^{-4}$, plateau-mean error $0.029\% \to 0.79\%$). The
curriculum therefore meets its narrow design objective without
improving the full-domain field or the extraction
(Sec.~\ref{sec:discussion}).
